\documentclass{article}
\usepackage{graphicx} % Required for inserting images
\usepackage{datetime} % Add this line to include the datetime package
\usepackage{amsmath}
\usepackage{amssymb}
\usepackage{amsthm}
\usepackage{tikz}
%\usepackage{hyperref}
\usepackage{xcolor} % Required for defining custom colors

\definecolor{darkblue}{rgb}{0.0, 0.0, 0.5} % Define a darker blue color

\usepackage[colorlinks=true, linkcolor=darkblue, urlcolor=darkblue, citecolor=darkblue]{hyperref} % Use the darker blue color

% Define theorem style
\theoremstyle{plain} % Bold title, italic body
\newtheorem{theorem}{Theorem}
\newtheorem{lemma}[theorem]{Lemma}
\newtheorem{corollary}[theorem]{Corollary}

\theoremstyle{definition} % Bold title, normal body
\newtheorem{definition}[theorem]{Definition}
\newtheorem{example}[theorem]{Example}


\title{Homework 7}
\author{Aaron Honculada}
\date{\monthname[\month] \the\year}

\begin{document}

\maketitle

\section*{Problem 1}
\subsection*{Determine the rules for addition and multiplication
of the congruence classes for the ring $\mathbb{Q}[x]/(x^2-2)$.}

We can represent the congruence classes as:
\[
    [a] = a + (x^2-2)
\]
where $a \in \mathbb{Q}[x]$.

To find the rules for addition and multiplication, we need to consider the representatives of the congruence classes.
Let's take two representatives:
\[
    [a] = a + (x^2-2)
\]
\[
    [b] = b + (x^2-2)
\]
We add and multiply these representatives to find the rules for addition and multiplication.

\begin{align*}
    [a] + [b] &= (a + (x^2-2)) + (b + (x^2-2)) \\
    &= (a+b) + (x^2-2) \\
    &= [a+b]
\end{align*}

\begin{align*}
    [a] \cdot [b] &= (a + (x^2-2)) \cdot (b + (x^2-2)) \\
    &= ab + a(x^2-2) + b(x^2-2) + (x^2-2)(x^2-2) \\
    &= [ab]
\end{align*}




\section*{Problem 2}
\subsection*{Show that $\mathbb{Q}[x]/(x^2)$ is not a field.}

\begin{proof}
    Consider the element $x$ in $\mathbb{Q}[x]/(x^2)$. We show that there is no element
    $y = a+bx \in \mathbb{Q}[x]/(x^2)$ such that $x \cdot y = 1$.
    
    For the sake of contradiction, assume that there exists such an element $y$.
    Then, we have:
    \[
        x \cdot (a+bx) = a \cdot x + b \cdot x^2 = a \cdot x + b \cdot 0 = a \cdot x
    \]
    This implies that:
    \[
        a \cdot x = 1
    \]
    This is a contradiction since the left hand side is a multiple of $x$ and the right hand side is a constant.
    Therefore, there is no non-zero element $y \in \mathbb{Q}[x]/(x^2)$ such that $x \cdot y = 1$.
    Therefore, $\mathbb{Q}[x]/(x^2)$ is not a field.
\end{proof}

\section*{Problem 3}
\subsection*{Explain why the class $[2x-3] \in \mathbb{Q}[x]/(x^2-2)$
is a unit and find its inverse.}

\begin{proof}
    We show there exists a polynomial $g(x) \in \mathbb{Q}[x]/(x^2-2)$ such that:
    \[
        (2x-3)g(x) = 1
    \]
    Using Euclidean algorithm, we find $(2x-3, x^2-2) = frac{1}{4} = 1 \pmod{x^2-2}$.
    Since the greatest common divisor is a constant then the class $[2x-3]$ is a unit.
    Given the linear combination $1 = 4(x^2-2) + (2x-3)(2x+3)$, we can express
    the inverse of $[2x-3]$ as:
    \[
        [2x-3]^{-1} = [2x+3] = [x+\frac{3}{2}].
    \]
\end{proof}

\subsection*{Explain why the class $[x^2+x+1] \in \mathbb{Z}_3[x]/(x^2+1)$
is a unit and find its inverse.}

\begin{proof}
    Using Euclidean algorithm, we are able to represent 1 as a linear combination of $x^2+x+1$ and $x^2+1$.
    \[
        1 = (x^2+x+1) - x(x+1)
    \]
    Therefore, the class $[x^2+x+1]$ is a unit and its inverse is $[x+1]$.
\end{proof}

\section*{Problem 4}
\subsection*{Show that $\mathbb{Q}(\sqrt{2})$ is isomorphic to $\mathbb{Q}[x]/(x^2-2)$.}

\begin{proof}
    We define a map $\phi: \mathbb{Q}[x]/(x^2-2) \to \mathbb{Q}(\sqrt{2})$ by:
    \[
        \phi(x) = \sqrt{2}
    \]
    We can verify that $\phi$ is a homomorphism.
    \begin{itemize}
        \item $\phi((a+bx) + (c+dx)) = \phi((a+c) + (b+d)x) = (a+c) + (b+d)\sqrt{2}$
            = $(a+b\sqrt{2}) + (c+d\sqrt{2}) = \phi(a+bx) + \phi(c+dx)$
        \item $\phi((a+bx)(c+dx)) = \phi((ac + 2bd) + (ad + bc)x) = (ac + 2bd) + (ad + bc)\sqrt{2}$
            = $(a+b\sqrt{2})(c+d\sqrt{2}) = \phi(a+bx)\phi(c+dx)$
    \end{itemize}

    We can verify that $\phi$ is injective.

    Suppose that $\phi(f(x)) = 0$. Then $f(x) \text{ maps to } 0$ in $\mathbb{Q}(\sqrt{2})$.
    This imples that $f(x)$ has a zero value when $x = \sqrt{2}$. Therefore, $f(x)$ must
    be a zero polynomial in $\mathbb{Q}[x]/(x^2-2)$ and therefore $f(x)$ is injective.
    
    For every element $a + b\sqrt{2} \in \mathbb{Q}(\sqrt{2})$, we can find a representative $a + bx \in \mathbb{Q}[x]/(x^2-2)$
    such that $\phi(a + bx) = a + b\sqrt{2}$. Therefore, $\phi$ is surjective.

    Since $\phi$ is a homomorphism, injective, and surjective, $\phi$ is an isomorphism.
    Therefore, $\mathbb{Q}[x]/(x^2-2) \cong \mathbb{Q}(\sqrt{2})$.
\end{proof}

\section*{Problem 5}
\subsection*{If $f(x) \in F[x]$ has degree $n$, prove that there exists
an extension field $E$ of $F$ such that $f(x)= c_0(x-c_1)(x-c_2)\cdots (x-c_n)$
for some (not necessarily distinct) $c_i \in E$. In other words, $E$ contains
all the roots of $f(x)$.}

\begin{proof}
    We know that for any field $F$, there exists an extension field $E$ of $F$, such 
    that non-constant polynomial $f(x) \in F[x]$ has a root in $E$.

    We can use induction to show that $f(x)$ has $n$ roots in $E$.

    Base case: $n = 1$.
    If $f(x)$ has degree 1, then $f(x) = ax + b$ for some $a, b \in F$.
    Since $f(x)$ has a root in $E$, we have $f(c) = 0$ for some $c \in E$.
    Therefore, $ac + b = 0$ and $c = -\frac{b}{a}$ and $f(x)$ has one root in $E$.

    Inductive step: Assume that for some $k \geq 1$, any non-constant polynomial of degree $k$ has $k$ roots in $E$.
    Consider a non-constant polynomial $f(x)$ of degree $k+1$.

    If $f(x)$ has a root in $E$, then we are done.
    Otherwise, we can use the division algorithm to write $f(x) = (x-c)g(x)$ for some $c \in E$ and $g(x) \in E[x]$.
    By the inductive hypothesis, $g(x)$ has $k$ roots in $E$.
    Therefore, $f(x)$ has $k+1$ roots in $E$.

    By induction, we have shown that any non-constant polynomial of degree $n$ has $n$ roots in $E$.

    We can conclude that $f(x)$ has $n$ roots in $E$ and can be written as $f(x) = c_0(x-c_1)(x-c_2)\cdots (x-c_n)$ for some (not necessarily distinct) $c_i \in E$.
\end{proof}

\section*{Problem 6}
\subsection*{List the disctinct principal ideals
in $\mathbb{Z}_5$}
\begin{itemize}
    \item $\left<0\right> = \{0\}$
    \item $\left<1\right> = \left<2\right> = \left<3\right> = \left<4\right> = \mathbb{Z}_5$
\end{itemize}

\subsection*{List the disctinct principal ideals
in $\mathbb{Z}_{12}$}
\begin{itemize}
    \item $\left<0\right> = \{0\}$
    \item $\left<1\right> = \left<5\right> = \left<7\right> = \left<11\right> = \mathbb{Z}_{12}$
    \item $\left<2\right> = \left<10\right> = \{0, 2, 4, 6, 8, 10\}$
    \item $\left<3\right> = \left<9\right> = \{0, 3, 6, 9\}$
    \item $\left<4\right> = \left<8\right> = \{0, 4, 8\}$
    \item $\left<6\right> = \{0, 6\}$
\end{itemize}
\section*{Problem 7}
\subsection*{If $I$ and $J$ are ideals in $R$, prove
that $I \cap J$ is an ideal in $R$.}

\begin{proof}
    Given that $I$ and $J$ are ideals in $R$, we know that both $I$ and $J$ are
    non-empty. 
    We show that $I \cap J$ is closed under subtraction. If $a, b \in I \cap J$, then $a, b \in I$ and $a, b \in J$.
    Since $I$ and $J$ are ideals, we have $a-b \in I$ and $a-b \in J$. Therefore, $a-b \in I \cap J$.

    We also show that $I \cap J$ has the absorption property. If $a \in I \cap J$ and $b \in R$, then $a \in I$ and $a \in J$.
    Since $I$ and $J$ are ideals, we have $ab \in I$ and $ab \in J$. Therefore, $ab \in I \cap J$.

    Since $I \cap J$ is closed under subtraction and has the absorption property, $I \cap J$ is an ideal of $R$.
\end{proof}

\section*{Problem 8}
\subsection*{Let $R$ be an integral domain and $a,b \in R$.
Show that $(a) = (b)$ if and only if $a=bu$ for some unit $u \in R$.}
\begin{proof}
    We first show that if $(a) = (b)$, then $a = bu$ for some unit $u \in R$.
    Since $(a) = (b)$, we have $a \in (b)$ and $b \in (a)$.
    Therefore, there exists $u \in R$ such that $a = bu$ and $b = au$.
    Since $R$ is an integral domain, we have $u$ is a unit.

    We now show that if $a = bu$ for some unit $u \in R$, then $(a) = (b)$.
    Since $a = bu$, we also have $b = au^{-1}$. Therefore, $a \in (b)$ and $b \in (a)$.
    Therefore, $(a) = (b)$.

    Therefore, $(a) = (b)$ if and only if $a = bu$ for some unit $u \in R$.
\end{proof}


\section*{Problem 9}
Let $J$ be the set of all polynomials with zero constant term
in $\mathbb{Z}[x]$.

\subsection*{Show that $J$ is the principal ideal $(x)$ in $\mathbb{Z}[x]$.}
\begin{proof}
    We define $J = \{f(x) \in \mathbb{Z}[x] \mid f(0) = 0\}$. This means that
    every polynomial in $J$ is divisible by $x$ and can be expressed as: $f(x) = xg(x)$ for some $g(x) \in \mathbb{Z}[x]$.

    Polynomials in $J$ are closed under subtraction. If $f(x), g(x) \in J$, then $f(x) - g(x) \in J$ since the
    constant term of $f(x) - g(x)$ is $0$.

    We also show that $J$ has the absorption property. If $f(x) \in J$ and $g(x) \in \mathbb{Z}[x]$, then $f(x)g(x) \in J$ since the
    constant term of $f(x)g(x)$ is still $0$.

    Next, we consider the ideal $(x)$. We can express any polynomial $f(x) \in \mathbb{Z}[x]$ as $f(x) = xg(x)$. 
    If $f(x) \in J$, then $f(0) = 0$ and $x$ divides $f(x)$. So there exists a $g(x) \in \mathbb{Z}[x]$ such that $f(x) = xg(x)$.
    Therefore, $J \subseteq (x)$.

    Conversely, any polynomial in the ideal $(x)$ of the form $xg(x)$ is in $J$.
    Therefore, $(x) \subseteq J$. Since both inclusions hold, we have $(x) = J$.

    Since $J$ is closed under subtraction and has the absorption property, $J$ is an ideal of $\mathbb{Z}[x]$.
    Furthermore, $J$ is principal since it is generated by $x$.
\end{proof}

\subsection*{Show that $\mathbb{Z}[x]/J$ consists of an infinite number
of distinct cosets, on for each number $n \in \mathbb{Z}$.}
\begin{proof}
    The elements of the quotient ring $\mathbb{Z}[x]/J$ are the cosets of $J$.
    These cosets are of the form $f(x) + J$ for some $f(x) \in \mathbb{Z}[x]$.
    Otherwise stated, polynomials $f(x) \cong g(x) \pmod{J} \iff f(0) = g(0)$ and
    each coset is uniquely determined by the constant term of the polynomial.

    Therefore there are an infinite number of distinct cosets, one for each integer $n \in \mathbb{Z}$.
    This also shows that $\mathbb{Z}[x]/J \cong \mathbb{Z}$.
\end{proof}

\section*{Problem 10}
Consider the ring of integers $\mathbb{Z}$.

\subsection*{Show that for each non-negative integer $n$,
$(n)$ is an ideal of $\mathbb{Z}$, and that all of these ideals are distinct.}

\begin{proof}
    Let $(n)$ be the set of all multiples of $n$. We know for that $n$, $(n)$ is
not empty. Furthermore, if $a, b \in (n)$, we can express $a$ and $b$ as $a = nk_1$ and $b = nk_2$ for some $k_1, k_2 \in \mathbb{Z}$.
Then, $a-b = nk_1 - nk_2 = n(k_1 - k_2) \in (n)$. We have shown that $(n)$ is closed under subtraction.
We also show that $(n)$ has the absorption property. If $a \in (n)$ and $b \in \mathbb{Z}$, then $ab = nk_1b = n(k_1b) \in (n)$.

We have shown that for each non-negative integer $n$, $(n)$ is an ideal of $\mathbb{Z}$.

We first consider the trivial case where $n = 0$. $(0)$ is unique in that it is
only contains the element $0$. Now we consider the two ideals $(n)$ and $(m)$ where $n, m > 0$.
We assume without loss of generality that $m < n$. Since $m < n$, we know $n \in (n)$, but $n \notin (m)$ if 
$m \nmid n$. Therefore, $(n) \neq (m)$. In the case where $n = m$, we have $(n) = (m)$ since there exists some
$k \in \mathbb{Z}$ such that $n = mk$. Therefore, all of the ideals are distinct.
\end{proof}

\subsection*{Prove that these are all of the ideals of $\mathbb{Z}$.}
\begin{proof}
    For each ideal $I$ of $\mathbb{Z}$, we show that if $s \in I$ then $(s) \subseteq I$.
    By the well-ordering principle, there must be a smallest positive element of $I$.
    Let $s$ be the smallest positive element of $I$. By definition, $(s)$ is an ideal of $\mathbb{Z}$ and
    contains all multiples of $s$. Since $s \in I$, we know that $(s) \subseteq I$.

    Let $a \in I$. We can express $a$ as $a = qs + r$ where $0 \leq r < s$. Since $a,s \in I$ and $I$ is closed
    under subtraction, we have: $r = a - qs \in I$. Since $s$ is the smallest positive element of $I$, we must have $r = 0$
    which implies that $a = qs \in (s)$. Therefore, every element $a \in I$ is 
    a multiple of $s$. The set of all multiples of $s$ is $(s)$. Therefore, $I \subseteq (s)$.

    Since $(s) \subseteq I$ and $I \subseteq (s)$, we have $(s) = I$. Therefore all ideals of $\mathbb{Z}$ of the form $(n)$, where $n \in \mathbb{Z}$ , they 
    are principal ideals.
\end{proof}


\end{document}